\section{The backward flow as a Bayesian inverse problem}\label{lo:sec:bayes}

\subsection{Positive linear inverse problem}

Discretise. Let $y_i=u_1(v_i)$ on a grid $v_i$ of spacing $h$ covering the data window, let $x_j\ge0$ be the unknown densities of $\mathfrak U_{\theta'}$ on a grid $w_j$ of the same spacing covering a larger support (cell masses $hx_j$), and let
\begin{equation}\label{lo:eq:plip}
y=Ax,\qquad A_{ij}=hP_{1-\theta'}(v_i-w_j).
\end{equation}
Every entry of $A$ is positive, and a column centred well inside the data window has $\sum_iA_{ij}\approx1$. This is a positive linear inverse problem in the sense of Vardi and Lee \cite{VardiLee93}: the same structure as emission tomography, image deblurring and the portfolio problem of their title, with the Cauchy kernel as the blur. Two features distinguish it from the usual instances. The data are exact, so the problem is one of existence rather than estimation. And the existence question is the zero free half plane $\Re s>\theta'$, the Riemann hypothesis when $\theta'=\tfrac12$: by Proposition~\ref{lo:prop:emres}(a), a positive solution of the continuous problem exists if and only if $\zeta$ has no zeros to the right of $\theta'$.

\subsection{Likelihood, prior, posterior}

The data $y_i$ are exact real numbers, so a Poisson model is a device, but a useful one. Read $hy_i$ as the expected count in cell $i$ of a Poisson process whose intensity is the Cauchy blur of an unknown positive measure: the complete data are counts $z_{ij}$ of mass originating in cell $j$ and observed in cell $i$, Poisson with means $hA_{ij}x_j$, and the observed counts are $\sum_jz_{ij}$. Replacing these counts by the data $hy_i$, the log likelihood of $x$ is, up to constants and the factor $h$,
\[
\ell(x)=\sum_i\bigl(y_i\log(Ax)_i-(Ax)_i\bigr),
\]
concave in $x$, and maximising it over $x\ge0$ minimises the generalised Kullback Leibler divergence $\mathrm{KL}(y\,\|\,Ax)=\sum_i[y_i\log(y_i/(Ax)_i)-y_i+(Ax)_i]$, which is defined for real data. The EM algorithm with the $z_{ij}$ as missing data is the Shepp and Vardi iteration \cite{ShVa82}
\begin{equation}\label{lo:eq:em}
x_j^{(k+1)}=\frac{x_j^{(k)}}{\sum_iA_{ij}}\sum_iA_{ij}\frac{y_i}{(Ax^{(k)})_i},
\end{equation}
which is also the Richardson and Lucy deblurring algorithm \cite{Richardson72,Lucy74}. It preserves positivity at every step, increases the likelihood monotonically, and, from a strictly positive start, converges to a maximum likelihood solution \cite{VSK85}. Csisz\'ar and Tusn\'ady \cite{CT84} identified it as alternating $I$ projection between the set of blurred positive measures and the set of distributions consistent with the data, an information theoretic analogue of von Neumann's alternating projection theorem for subspaces of a Hilbert space.

A Bayesian completes the model with a prior on $x$, the conjugate choice being a product of gamma laws. A sampler needs integer counts, which exact data do not supply: fix an exposure $\tau>0$ and let $N_i$ be the integer nearest to $\tau hy_i$, modelled as $N_i=\sum_jZ_{ij}$ with independent $Z_{ij}$, Poisson with means $\tau hA_{ij}x_j$. Given $x$ and $N$, the vector $(Z_{ij})_j$ is multinomial with total $N_i$ and probabilities $A_{ij}x_j/(Ax)_i$; given $Z$, with independent $\mathrm{Gamma}(\alpha_j,\beta_j)$ priors, $x_j$ is gamma with shape $\alpha_j+\sum_iZ_{ij}$ and rate $\beta_j+\tau h\sum_iA_{ij}$. Alternating the two steps is a data augmentation Gibbs sampler, the Monte Carlo counterpart of \eqref{lo:eq:em} in which the E step is replaced by simulation: the deterministic EM fixed point becomes a posterior over positive measures, and the uncertainty that EM hides becomes the spread of the posterior. That spread depends on the exposure $\tau$, which exact data do not determine. The sampler is a method; no run of it is reported here (Problem~\ref{lo:prob:sampler}).

The generative Bayesian computation of the next section amortises the posterior of a different model. Instead of running a Gibbs chain for each data set, one simulates pairs from a prior on configurations of zeros and the forward map, trains a quantile network to represent conditional quantiles of chosen functionals given the data, and evaluates it at the observed data. The three methods share the forward map but not the prior or the unknown: EM gives the posterior mode under a flat prior on the grid masses, data augmentation samples the posterior under a gamma prior on the same masses, and GBC regresses functionals of a GUE configuration of zeros on simulated data.

\subsection{What positivity buys}

The Fourier backward flow of Section~\ref{lo:sec:backward} ignores positivity, which is why the sharp cutoff rings below zero at $\theta'=\tfrac12$ and why its error is amplified by $e^{(1-\theta')\Omega}$. The EM inversion enforces positivity and uses the whole frequency content of the data, so it is the natural positivity preserving backward flow, and its total residual is the diagnostic of Proposition~\ref{lo:prop:emres}.

\begin{proposition}[the EM residual as a diagnostic]\label{lo:prop:emres}
Let $\theta'\in[\tfrac12,1)$ and $a=1-\theta'$.
\begin{enumerate}
\item[(a)] There is a positive measure $x$ on $\R$ with $\int x(dv)/(1+v^2)<\infty$ and $P_a*x=u_1$ if and only if $\zeta$ has no zeros in $\Re s>\theta'$. It is then unique and equals $\mathfrak U_{\theta'}$.
\item[(b)] Let $h_k\to0$ and $L_k\to\infty$ with $h_k^2L_k\log L_k\to0$. For positive measures $x=\sum_jx_j\delta_{jh_k}$ on the grid $h_k\mathbb Z$ with $\sum_jx_j/(1+j^2h_k^2)<\infty$, let $m_i(x)=(P_a*x)(v_i)$ at the data points $v_i\in h_k\mathbb Z\cap[-L_k,L_k]$, let $y_i=u_1(v_i)$, and let
\[
D_k=\inf_x\,h_k\sum_i\Bigl(y_i\log\frac{y_i}{m_i(x)}-y_i+m_i(x)\Bigr)
\]
be the total, not normalised, Kullback Leibler residual. Then $D_k\to0$ if $\zeta$ has no zeros in $\Re s>\theta'$, and $\liminf_kD_k>0$ if it has one.
\end{enumerate}
\end{proposition}

\begin{proof}
(a) is Theorem~\ref{tf:thm:deconv}(a) with $\varepsilon_0=\tfrac12$ and $\varepsilon_1=\theta'-\tfrac12$.

Write $d(y,m)=y\log(y/m)-y+m$; then $(\sqrt y-\sqrt m)^2\le d(y,m)\le(y-m)^2/m$ for $y,m>0$. Also $|\partial_t\log P_a(t)|=2|t|/(a^2+t^2)\le1/a$, so $P_a(t')$ lies between $P_a(t)e^{-|t-t'|/a}$ and $P_a(t)e^{|t-t'|/a}$.

(b), no zeros in $\Re s>\theta'$. Let $x_j=\mathfrak U_{\theta'}([jh_k-h_k/2,jh_k+h_k/2))$, which satisfies the summability condition because $\mathfrak U_{\theta'}$ does. By (a), $y_i=\int P_a(v_i-w)\,\mathfrak U_{\theta'}(dw)$, and moving each $w$ to the centre of its cell changes $P_a(v_i-w)$ by a factor in $[e^{-h_k/2a},e^{h_k/2a}]$, so $|m_i(x)-y_i|\le(e^{h_k/2a}-1)y_i$ and $d(y_i,m_i(x))\le Ch_k^2y_i$ for $h_k\le a$. As $u_1=P_a*\mathfrak U_{\theta'}$ gives $|u_1'|\le u_1/a$, we have $h_k\sum_iy_i\le e^{h_k/2a}\int_{-L_k-1}^{L_k+1}u_1\,dv$, which is $O(L_k\log L_k)$ by the Riemann von Mangoldt formula: the zeros with $|\gamma_\rho|\le2L_k+2$ contribute at most their number, and the others $O(\log L_k)$. Hence $D_k=O(h_k^2L_k\log L_k)\to0$. Atoms of $\mathfrak U_{\theta'}$, at zeros with $\sigma_\rho=\theta'$, are allowed.

(b), a zero in $\Re s>\theta'$. Suppose that $D_k\to0$ along a subsequence, and choose grid measures $x_k$ with total residual at most $D_k+1/k$. The function $F_k=P_a*x_k$ is smooth and positive with $|F_k'|\le F_k/a$, and $h_k\sum_{|v_i|\le R}(\sqrt{y_i}-\sqrt{F_k(v_i)})^2\to0$ for every $R>0$. Since $F_k\ge e^{-1/a}F_k(0)$ on $[-1,1]$, the case $R=1$ shows that $F_k(0)=(a/\pi)\int x_k(dw)/(a^2+w^2)$ is bounded. So the measures $\nu_k(dw)=x_k(dw)/(a^2+w^2)$ have bounded mass, and along a further subsequence they converge weakly on the compactification $\R\cup\{\infty\}$ to a finite measure $\nu$. Now $F_k(v)=\int\kappa_v\,d\nu_k$ with $\kappa_v(w)=(a/\pi)(a^2+w^2)/(a^2+(v-w)^2)$, which is jointly continuous on $[-R,R]\times(\R\cup\{\infty\})$ with $\kappa_v(\infty)=a/\pi$. Hence $F_k\to F$ uniformly on $[-R,R]$, where $F=P_a*x+c$ with $x=(a^2+w^2)\nu$ on $\R$ and $c=(a/\pi)\nu(\{\infty\})\ge0$, and Riemann sums give $\int_{-R}^R(\sqrt{u_1}-\sqrt F)^2\,dv=0$. As $R$ is arbitrary and $P_a*(c\,dv)=c$, the positive measure $x+c\,dv$ has $\int(x+c\,dv)/(1+v^2)<\infty$ and $P_a*(x+c\,dv)=u_1$, which (a) excludes. So $\liminf_kD_k>0$.
\end{proof}

The total residual cannot be replaced by the residual per unit length of the data window: when only finitely many zeros lie to the right of $\theta'$, the latter tends to zero, since an atom of mass $m_\rho$ at each of their ordinates leaves the defect $\sum_\rho(P_{1-\sigma_\rho}-P_a)(v-\gamma_\rho)$, whose total residual is finite. The proof of (b) also shows that the residual restricted to some fixed window $[-R,R]$ stays bounded away from zero, and the same argument works for every $R>0$ since $u_1$ and $P_a*x+c$ are real analytic, but the proposition gives no rate. The runs below, with a fixed data window and a finite unknown support, illustrate it; no finite computation certifies either alternative.

\subsection{Numerical EM inversion}\label{lo:sec:emnum}

{\emergencystretch=8em The runs are \texttt{code/lo\_\allowbreak em.py} (\texttt{data/lo\_\allowbreak em\_\allowbreak table.json}, \texttt{data/lo\_\allowbreak em\_\allowbreak inject.json}), with the EM iteration \eqref{lo:eq:em} of \texttt{code/cx\_em.py} (the convolution by FFT agrees with the dense product to $2\times10^{-15}$). Data are $u_1$ on $[100,300]$, the unknown is supported on $[80,320]$, spacing $0.05$, flat start with the total mass matched. The residual printed is the plain sum $\sum_i[y_i\log(y_i/(Ax)_i)-y_i+(Ax)_i]$ over the $4001$ data points, that is $\mathrm{KL}(y\,\|\,Ax)$ of Proposition~\ref{lo:prop:emres} divided by the spacing $h$, not per point and not per unit length; it changes by $1.6$ per cent between spacings $0.05$ and $0.1$, so two digits at most are meaningful. A peak is a local maximum of prominence at least $1$ per cent of the maximum, and $\min x$ is the smallest $x_j$.\par}

\begin{table}[ht]
\centering
\small
\begin{tabular}{cccccc}
\toprule
$\theta'$ & steps & KL residual & peaks & max error & $\min x$\\
\midrule
0.75 & 300 & $2.4\times10^{-5}$ & 109 of 109 & 0.040 & $7.7\times10^{-2}$\\
0.50 & 300 & $4.9\times10^{-2}$ & 109 of 109 & 0.085 & $2.7\times10^{-5}$\\
0.50 & 3000 & $4.0\times10^{-3}$ & 109 of 109 & 0.030 & $9.9\times10^{-12}$\\
\bottomrule
\end{tabular}
\caption{EM inversion of $u_1$ on $[100,300]$ (computed). The source printed $4.1\times10^{-3}$ in the last row; the recomputation gives $4.05\times10^{-3}$, and the grid sensitivity of the residual is of that size.}\label{lo:tab:em}
\end{table}

At $\theta'=0.75$ the iteration converges quickly to a strictly positive density with one peak per zero; at $\theta'=0.5$ the target is atomic, the iteration sharpens slowly, as EM does when the truth lies on the boundary of the positive cone, and after $3000$ steps the mass has concentrated at the zeros to $0.03$ with the density between them driven to $10^{-11}$. There is no ringing.

With an off line pair injected at $\sigma=0.75$ and $\theta'=0.75$, the source reports a residual after $3000$ steps of $1.6\times10^{-5}$ against $2.3\times10^{-7}$ for the real $u_1$, a factor of about $70$; the recomputation gives $1.62\times10^{-5}$ and $2.27\times10^{-7}$ (ratio $71$, and $53$ on the data window $[150,250]$) for a quartet at $\sigma=\theta'$ that keeps the total mass, and the source's local misfit near the injected zero could not be reproduced, because its definition is not given. This zero lies on the level, not to its right, so a positive solution still exists: at level one the pair adds $P_{1/4}(v-\gamma)+P_{3/4}(v-\gamma)$, which is $P_{1/4}*(\delta_\gamma+P_{1/2}(\cdot-\gamma))$. The larger residual therefore measures the slow concentration of EM onto an atom, as in the $\theta'=0.5$ rows, and does not illustrate Proposition~\ref{lo:prop:emres}(b), which needs $\sigma>\theta'$.

The test was therefore redesigned. The two consecutive line zeros nearest the centre of the window ($\gamma=198.0153$ and $201.2648$, midpoint $199.6400$) are removed and replaced by the quartet $\rho,1-\bar\rho,\bar\rho,1-\rho$ with real part $\sigma$ at the midpoint (widths $1-\sigma$ and $\sigma$ at the midpoint and at its mirror image, total mass unchanged to $10^{-4}$). The control puts two line zeros at the midpoint. Table~\ref{lo:tab:inject} gives the total residual at $10000$ steps and the ratio of quartet to control at $300$, $1000$, $3000$ and $10000$ steps, on the data windows $[100,300]$ (unknown $[80,320]$) and $[150,250]$ (unknown $[130,270]$).

\begin{table}[ht]
\centering
\footnotesize
\setlength{\tabcolsep}{3pt}
\begin{tabular}{cccccc}
\toprule
$\theta'$, $\sigma$ & window & real $u_1$ & control & quartet & ratio at $300,\dots,10^4$ steps\\
\midrule
0.50, 0.75 & $[100,300]$ & $1.09\times10^{-3}$ & $1.08\times10^{-3}$ & $1.85\times10^{-2}$ & 1.4, 2.4, 5.4, 17\\
0.50, 0.75 & $[150,250]$ & $5.0\times10^{-4}$ & $5.0\times10^{-4}$ & $1.78\times10^{-2}$ & 1.8, 3.8, 10, 36\\
0.75, 0.85 & $[100,300]$ & $3.2\times10^{-8}$ & $3.3\times10^{-8}$ & $7.03\times10^{-3}$ & 366, 4000, 31000, 220000\\
0.75, 0.85 & $[150,250]$ & $5.9\times10^{-8}$ & $5.9\times10^{-8}$ & $7.02\times10^{-3}$ & 627, 4500, 23000, 120000\\
0.75, 0.75 & $[100,300]$ & $3.2\times10^{-8}$ & $3.3\times10^{-8}$ & $4.4\times10^{-6}$ & 9.0, 30, 72, 136\\
\bottomrule
\end{tabular}
\caption{Quartet injection test with controls (computed, \texttt{code/lo\_em.py inject}); total residual after $10000$ EM steps. The last row is the boundary case $\sigma=\theta'$.}\label{lo:tab:inject}
\end{table}

\begin{resultpar}\label{lo:res:inject}The redesigned test separates the situations. At $\theta'=0.75$ and $\sigma=0.85>\theta'$ the quartet residual stops at $7.03\times10^{-3}$, the same for both windows to $0.3$ per cent: a floor, with no positive solution, as Proposition~\ref{lo:prop:emres}(b) predicts for a zero to the right of $\theta'$, while the control and the real data decay below $10^{-7}$. At $\theta'=0.5$ and $\sigma=0.75$ the control falls like $1/k$ ($4.9\times10^{-2}$, $1.3\times10^{-2}$, $4.0\times10^{-3}$, $1.1\times10^{-3}$ at $300$, $1000$, $3000$, $10000$ steps) while the quartet leaves $6.9\times10^{-2}$, $3.1\times10^{-2}$, $2.2\times10^{-2}$, $1.85\times10^{-2}$; the floor is not reached within $10000$ steps and the ratio keeps growing, so this row shows slower decay than the control and not a plateau. The boundary case $\sigma=\theta'$ keeps decaying, like $1/k$, but more slowly than the control (ratio $9.0$ at $300$ steps, $136$ at $10000$), because a positive solution exists and EM is slow onto an atom. The plain sum is $1/h$ times the total residual of the proposition, a fixed factor at fixed spacing; a finite run certifies neither alternative.\end{resultpar}

\section{Generative Bayesian computation for the backward flow}\label{lo:sec:gbc}

The backward flow is a deconvolution, and a deconvolution is a Bayesian inverse problem once one says what is believed about the unknown. The unknown is the configuration $Z=\{\gamma\}$ of zeros in a window; the forward map is the Cauchy semigroup,
\[
Z\mapsto\sum_{\gamma\in Z}\mathrm{Cauchy}(\gamma,\tfrac12)
\]
on the window, plus the contribution of the zeros outside the window, and nothing else: by Theorem~\ref{lo:thm:density} the zero sum is the whole of $u_1$, and the smooth part $g_1$ belongs to the prime side decomposition \eqref{lo:eq:u1}. The data are the values of $u_1$ on a grid, exact to rounding. A prior for $Z$ is supplied by Montgomery's pair correlation \cite{Montgomery73} and Odlyzko's computations \cite{Odlyzko87}: locally the zeros look like the eigenvalues of a large GUE matrix, unfolded to density $\frac1{2\pi}\log\frac{v}{2\pi}$.

Generative Bayesian computation (Polson and Sokolov \cite{PS23,PSGBC26}) handles such a problem without likelihoods or Markov chains. One simulates $(Z_i,y_i)$ from the prior and the forward map, trains a deep quantile network $q_\tau(y)$ to predict functionals of $Z$ from $y$ at quantile levels $\tau$, and evaluates the trained network at the observed $y$. The network is meant as an amortised posterior: by the noise outsourcing theorem \cite{Kallenberg02} any conditional law is the push forward of a uniform variable through a measurable map, and the quantile network learns that map for the functionals of interest; in practice it estimates marginal conditional quantiles of those functionals, not their joint law. Here the functionals are the counts of zeros in eight unit bins at the centre of a window of length $16$, which determine the zero positions to within the bin width and are what a deconvolution has to get right. The formulation is heuristic; what follows is a computation.

\subsection{Implementation}

The run is \texttt{code/lo\_gbc.py} (\texttt{data/lo\_gbc.json}), with two seeds, in $170$ seconds on $8$ cores. Training set: $24{,}000$ windows. For each window a density $d$ is drawn uniformly from $[0.45,0.8]$, which covers the heights $150\le v\le350$, where the zero density $(2\pi)^{-1}\log(v/2\pi)$ runs from $0.50$ to $0.64$ (and in fact the heights $106$ to $958$); a GUE matrix $H=(A+A^*)/2$ of size $M=400$, with $A=X+iY$ and $X,Y$ independent standard normal matrices, so that the off diagonal entries have variance $1$ and the spectral edge is $2\sqrt M=40$; the eigenvalues $x$ are unfolded to unit spacing by the semicircle law, $e=MF(x/(2\sqrt M))$ with $F$ the semicircle distribution function, and rescaled to density $d$ (the measured number of zeros per unit length divided by $d$ is $0.998$ over the $24{,}000$ windows); the forward map is evaluated on the grid of spacing $0.1$ as the sum of $P_{1/2}(t-z)$ over the simulated zeros within $150/d$ of the centre, plus the Cauchy smoothing of the mean density $d$ for the zeros beyond; the bin counts are the targets. Dividing the eigenvalues by $\sqrt{2M}$ instead of $2\sqrt M$ would give the density $d/\sqrt2$, which does not cover the real density above $v=220$. Network: two hidden layers of $256$ rectified units, $24$ outputs (eight bins at $\tau=0.1,0.5,0.9$), pinball loss, Adam, $25$ epochs, plain numpy. Test set: windows of the real $u_1$ centred at $v=150,158,\dots,342$ ($25$ windows, $200$ bins), with the true counts from the zeros. The forward map with the real zeros within $150/d$, $d=(2\pi)^{-1}\log(c/2\pi)$ at the window centre $c$, and the mean density beyond differs from $u_1$ by at most $4.4\times10^{-4}$ on these windows (mean $1.9\times10^{-4}$, against $\max u_1=1.07$): the modelling error of the outside zeros, the mirror zeros and the density gradient.

\subsection{Results}

\begin{resultpar}\label{lo:res:gbc}The two seeds are given as seed 1 and seed 2. The pinball loss at epoch $25$ is $0.0157$ and $0.0166$ on the training simulations, and $0.0161$ and $0.0175$ on held out simulations. On the $200$ real bins the rounded posterior median equals the true count in $196$ and $197$ bins ($98.0$ and $98.5$ per cent) and is within one of it in all $200$; the errors are $3$ undercounts and $1$ overcount, and $3$ undercounts, and about half of them fall in the five bins that hold two or more zeros, where two Cauchy bumps of width $\tfrac12$ merge into one. On held out simulations the quantile levels are calibrated to within about $0.05$ for integer counts: $P(N\le q_{0.1})=0.086$ and $0.095$, $P(N\le q_{0.5})=0.486$ and $0.449$, $P(N\le q_{0.9})=0.892$ and $0.891$, and the interval $[q_{0.1},q_{0.9}]$ covers $0.806$ and $0.796$ (nominal $0.8$); on the real windows it covers $0.805$ and $0.775$. The first run of the source reported $68$ per cent exact bins, $99$ per cent within one, and uncalibrated quantiles, with training inputs and test inputs that did not match (the smooth part $g_1$ was counted twice); with the corrected model the network counts almost exactly, and the earlier numbers are not those of this model.\end{resultpar}

A misfit statistic, the root mean square distance over $|t|\le4$ between $y$ and the forward map of the reconstruction that places the rounded median counts at the bin centres, is $0.173$ (standard deviation $0.031$) on the real windows and $0.331$ ($0.154$) and $0.307$ ($0.132$) on the same windows with an injected quartet at $\sigma=0.75$ (total mass kept), for the two seeds; with two line zeros placed at the midpoint instead it is $0.190$ and $0.186$. The injected misfit exceeds the real one in $20$ and $21$ of the $25$ windows. The source's statistic, which did not separate real from injected windows ($0.31$ against $0.34$), is not defined in the source and was not reproduced; the statistic here is a plug in misfit and not a posterior predictive check proper.

This is a first implementation that counts the zeros in unit bins and does not locate them. The EM inversion of Section~\ref{lo:sec:emnum}, which uses the same positivity and no prior, locates the zeros to $0.03$, while GBC gives bin counts only, which throw away the fine structure of the bumps, and the two are not compared on one metric. A GBC inversion that targets the EM solution itself, or the positions of the zeros as continuous outputs, with the EM reconstruction supplied as a feature, is the right next step and is left as an open problem.

\subsection{What the Bayesian view adds}

The GUE prior encodes the critical line: every simulated configuration has all its zeros at real part $\tfrac12$, so every simulated $y$ is a sum of Cauchy bumps of width exactly $\tfrac12$. A zero at $\sigma\ne\tfrac12$ produces, with its mirror, bumps of widths $1-\sigma$ and $\sigma$ in place of the two bumps of width $\tfrac12$ of two zeros on the line. Such a window is atypical under the prior, and a misfit statistic, the distance between the data and the forward map of the posterior median configuration, is meant to flag it; in the run above it separates the injected windows from the real ones on average, and a posterior predictive check proper would simulate from the posterior rather than plug in its median. The method can do two things that the Fourier deconvolution does not: regularise the ill posed inversion with a prior rather than with a frequency cutoff, and report uncertainty and misfit, to the extent that its quantiles are calibrated. What it cannot do is establish the absence of off line zeros. The prior assumes them away, and the misfit statistic can only flag a window where the data disagree with that assumption, at the resolution the window allows.

In the language of the book on Hilbert's sixth problem \cite{PSH626}, the forward map is the Poisson semigroup, the subordination of the heat semigroup by the stable one half law, and $\mathrm{Cauchy}(\gamma,a)$ is the law of $\gamma+B(T_a)$ for a Brownian motion $B$ and the hitting time $T_a$ of level $a$ by an independent Brownian motion. Under RH the level $\theta$ Thorin measure is thus the zero measure run through Brownian motion to the hitting time of level $\theta-\tfrac12$, and the backward flow is the unstopping of that Brownian motion; without RH the zeros carry widths $\theta-\sigma_\rho$ and the unconditional version is Section~\ref{bp:sec:planar}. That book reads this in thermodynamic terms; for the infinite measures here entropy is not defined, and the reading is heuristic, not an independent handle on RH.

\section{Is this a path to RH?}\label{lo:sec:path}

No, and it is worth being precise about why, because the chapter produces several objects that look like handles.

The level one measure is explicit and unconditional. Its positivity, which is what the Thorin theory gives for free, is automatic: by the Euler product every zero has $\sigma_\rho\le1$, so every term of \eqref{lo:eq:Utheta} at $\theta=1$ is a positive measure. The nonvanishing of $\zeta$ on the line $\Re s=1$ is equivalent to something else, the absence of atoms: $\mathfrak U_1$ has no atoms, and then has a continuous density, if and only if $\zeta(1+it)\ne0$ for every real $t$. Indeed a zero $1+i\gamma_0$ would contribute the atom $\mathrm{Cauchy}(\gamma_0,0)=\delta_{\gamma_0}$, which the other terms, all positive, cannot cancel, and if every width $1-\sigma_\rho$ is positive the proof of Theorem~\ref{lo:thm:density} gives a continuous density. Everything to the left of the line enters only through the backward flow.

The backward flow is a deconvolution by $e^{(1-\theta')|\omega|}$. Its convergence at level $\theta'$ is equivalent, by Theorem~\ref{lo:thm:dichotomy}, to the zero free half plane $\Re s>\theta'$, and the proof of that equivalence is the pole argument of Chapter~\ref{chap:lv}. On the prime side the flow is the passage from the weights $\Lambda(n)/n$ to $\Lambda(n)n^{-\theta'}$, and its convergence is equivalent to the classical bound $\psi(x)-x=O(x^{\theta'+\varepsilon})$ for every $\varepsilon>0$. The dictionary is exact and nothing is lost, but nothing is gained either: the error term in the prime number theorem has become the growth of a Fourier multiplier.

The Bayesian inversion, whether by EM, by data augmentation, or by generative Bayesian computation, is a numerical method for the inverse problem. Its residual is a diagnostic equivalent to the zero free half plane (Proposition~\ref{lo:prop:emres}), but equivalent is the operative word: the total residual stays bounded away from zero as the discretisation is refined exactly when there is a zero to the right of $\theta'$, and showing that it tends to zero at every $\theta'>\tfrac12$ is RH. The GUE prior of the GBC version assumes the critical line outright. These are tools for computation and verification, not for proof.

What the chapter does contribute is narrower: the prime formula of Propositions~\ref{bp:prop:thorin} and~\ref{bp:prop:sin2} extended to every level $\theta\ge1$ and to $v=0$; the explicit formula of Theorem~\ref{dc:thm:weil} read at level one, for one analytic test function; a dichotomy theorem for the backward flow; the total residual theorem for the discretised inversion; and a computation in which $109$ zeros are read off the line $\Re s=1$ by deconvolution to $0.03$, with a first GBC implementation that recovers bin counts. The question that remains is Problem~\ref{lv:prob:down}, now sharpened: find a property of the explicit measure \eqref{lo:eq:u1}, visible from the primes, that is preserved by the backward Poisson flow and forces it to remain a measure. Positivity is not such a property, since it is not preserved backward. We do not know one.

\section{Open problems}\label{lo:sec:open}

\begin{problem}\label{lo:prob:sampler}
Implement the data augmentation sampler of Section~\ref{lo:sec:bayes} for the real $u_1$ and compare its posterior spread at $\theta'=0.75$ and $\theta'=0.5$ with the EM point estimate; then build the GBC inversion with the EM reconstruction as an input feature and zero positions as continuous outputs, and calibrate its quantiles on held out simulations.
\end{problem}

\begin{problem}\label{lo:prob:selberg}
Characterise, without reference to zeros, the functions of the form \eqref{lo:eq:u1}: which densities $\frac1\pi[g(v)-\lim_X(\sum_{n\le X}a_n\cos(v\log n)-\sin(v\log X)/v)]$ with $a_n\ge0$ supported on prime powers are Thorin densities at level one of a function $\Xi(1)/\Xi(1+\sqrt s)$ with $\Xi$ in the Selberg class? The Thorin conditions constrain the growth and the integrability near $u=0$ but say nothing about the arithmetic of the $a_n$. (The wording is the author's; the Selberg class needs the completed function, the gamma factors, the conductor and the poles to be specified, and the sign of the coefficients of $-L'/L$.)
\end{problem}

\begin{problem}\label{lo:prob:rate}
Quantify the backward flow. For $\theta'\in(\theta_0,1)$ give the rate of locally uniform convergence of $B_{\theta',\Omega}\mathfrak U_1$ to $\mathfrak U_{\theta'}$ in terms of $\Omega$ and $\theta'-\theta_0$, and give the matching lower bound in case (ii): with $\sigma_{\max}=\sup\sigma_\rho>\theta'$ attained by a zero, show that $\sup_{|v|\le V}|B_{\theta',\Omega}\mathfrak U_1(v)|$ grows like $e^{(\sigma_{\max}-\theta')\Omega}$ up to polynomial factors. For one zero with $d=\sigma-\theta'>0$ the Fej\'er inverse at distance $t$ from its ordinate is $(\pi\Omega)^{-1}\Re[(e^{z\Omega}-1-z\Omega)/z^2]$ with $z=d+it$, which at $t=0$ is asymptotic to $e^{d\Omega}/(\pi d^2\Omega)$; with infinitely many zeros the supremum need not be attained and the cancellations among them are part of the problem.
\end{problem}

\begin{problem}\label{lo:prob:delta}
Replace the GUE prior by a prior that allows a fraction of zeros off the line, for instance pairs at $\tfrac12\pm\delta+i\gamma$ with $\delta$ drawn from a prior on $[0,\tfrac12)$, and train the quantile network to report the posterior of $\delta$ for each detected zero. The posterior of $\delta$ at the real $u_1$ is then a window by window numerical diagnostic at the resolution the data allow; being conditional on the prior, it is not a certificate.
\end{problem}
